\documentclass[]{article}

\usepackage{amsmath}
\usepackage{multirow}
\usepackage{natbib}
\usepackage{graphicx} 


\begin{document}

In this paper, we addressed the information loss inherent in traditional summary statistics, such as the two-point correlation function, by directly modeling the three-dimensional positions of galaxies. We introduced a framework that treats the galaxy distribution as a mass-conditioned spatial point process, allowing for a more detailed probe of the internal structure of dark matter halos.

Our analysis was performed on a suite of ten synthetic galaxy catalogs where the underlying Halo Occupation Distribution (HOD) parameters were known. We employed a Neyman-Scott process model, where the mean number of satellite galaxies follows a power-law with host halo mass and their radial positions are described by an exponential profile. Using Maximum Likelihood Estimation, we first demonstrated that our method successfully recovers the input HOD parameters governing the satellite population, albeit with a small, systematic bias that is well-understood as an effect of modeling a discrete process over a continuous mass function. This result served to validate our likelihood-based approach.

The primary findings of this work reveal important, mass-dependent features of halo structure. First, through a formal model comparison using the Akaike Information Criterion, we found decisive statistical evidence ($\Delta \text{AIC} = 136.13$) that the radial concentration of satellite galaxies increases with the mass of their host halo. This result challenges the common assumption of self-similarity in the spatial distribution of satellite galaxies, suggesting that physical processes shaping satellite orbits vary across the halo mass spectrum. Second, by employing a marked correlation function with luminosity as the mark, we quantified the spatial segregation of galaxies within their halos. The analysis showed that more luminous galaxies are preferentially located closer to the centers of their host halos. Finally, a residual analysis in the radial range of 5-10 Mpc/h, a region excluded from our model fit, precisely quantified the breakdown of the 1-halo model. The severe underprediction of galaxy counts at these scales highlights the transition to the 2-halo regime, where clustering between neighboring halos becomes the dominant contribution.

We have learned that direct likelihood-based modeling of spatial point patterns is a powerful and sensitive alternative to methods based on summary statistics. This framework can extract detailed astrophysical information about the mass-dependent properties of galaxy halos directly from catalog-level data. The evidence for a break in the self-similarity of satellite profiles demonstrates the potential of this approach to test foundational assumptions in galaxy formation models. This work establishes our method as a robust tool for analyzing the rich datasets from next-generation cosmological surveys, enabling more precise constraints on the galaxy-halo connection.

\end{document}
                